\documentclass[12pt,a4paper]{article}
\usepackage[utf8]{inputenc}
\usepackage[spanish]{babel}
\usepackage{geometry}
\usepackage{enumitem}
\usepackage{xcolor}

\geometry{a4paper, margin=1in}

\title{Comunicado Oficial: Procedimiento de Entrega de Tareas \\ \large Curso: Seguridad Informática}
\author{Prof. César Rodríguez}
\date{\today}

\begin{document}

\maketitle

\section*{Aviso Importante sobre la Entrega de Módulos}

Estimados estudiantes,

Me he percatado, a través de mensajes en nuestro grupo de WhatsApp, que varios de ustedes han editado y guardado correctamente los scripts de Python asignados a sus grupos en sus equipos locales, pero \textbf{no están realizando el Pull Request (PR) en GitHub}.

Quiero recordarles que notificar por WhatsApp que han terminado su código \textbf{no constituye una entrega formal} y no permite que sus cambios se integren a la Suite de Auditoría. Guardar el archivo en su computadora no actualiza el repositorio central.

\section*{Procedimiento Obligatorio de Entrega}

Para que su trabajo sea evaluado e integrado al orquestador (\texttt{auditoria.py}), es estrictamente necesario seguir el flujo de control de versiones. Una vez que su código funcione localmente, deben ejecutar en su terminal:

\begin{enumerate}
    \item \textbf{Añadir los cambios:} \texttt{git add modulos/su\_archivo.py}
    \item \textbf{Guardar en el historial:} \texttt{git commit -m 'Descripción de los cambios'}
    \item \textbf{Subir a GitHub:} \texttt{git push origin su\_rama}
    \item \textbf{Crear el Pull Request:} Ir a GitHub y presionar el botón verde 'Compare \& pull request'.
\end{enumerate}

Alternativamente, para simplificar el proceso, recuerden que pueden utilizar el script de automatización provisto en la raíz del proyecto:
\begin{verbatim}
python3 git_update.py 'Mensaje descriptivo de su trabajo'
\end{verbatim}

\section*{Consecuencias de no realizar el Pull Request}

\begin{itemize}[label=$\bullet$]
    \item Si el código no está en un PR, \textbf{no existe} para efectos de calificación.
    \item Retrasan el trabajo de los demás grupos y la integración de las fases.
    \item WhatsApp es un canal de comunicación para consultas rápidas, no una plataforma de control de versiones ni de recepción de tareas.
\end{itemize}

Por favor, revisen el documento \texttt{docs/guia\_pull\_requests.pdf} si tienen dudas con el proceso paso a paso. Quedo atento a la recepción formal de sus trabajos directamente en GitHub.

\vspace{1cm}
\noindent Atentamente, \\
\vspace{0.5cm} \\
\textbf{Prof. César Rodríguez}

\end{document}